\chapter{Introduction}
\label{ch:intro}

\section{Brief}
Air pollution is the environmental risk factor that shortens the most lives, and
South Asia carries a disproportionate share of that burden. In Pakistan, the
winter months bring a persistent smog season in which the Air Quality Index (AQI)
in the large cities of the Punjab plain routinely reaches values that public
health agencies classify as unhealthy or worse for days at a time. The people
most affected by it - children, the elderly, and anyone with a respiratory or
cardiac condition - have very little information with which to protect
themselves. What information does exist is almost invariably a \emph{nowcast}: a
single number describing the air as it is at this moment, at one monitoring
location, offered without explanation and without any statement of what the air
will be like tomorrow.

That is the wrong shape of information for the decisions people actually make.
Whether to send a child to school on foot, whether to move an outdoor shift
indoors, whether an asthmatic should carry a mask or reschedule an errand: these
are decisions taken hours or days in advance, and they need a forecast, not a
reading. A number that tells someone the air is dangerous \emph{now} arrives too
late to be acted upon.

This thesis presents the \projname, an end-to-end machine-learning system that
forecasts the United States Environmental Protection Agency AQI up to 72 hours
ahead for 22 cities spanning every province of Pakistan together with the
Islamabad Capital Territory, Gilgit-Baltistan and Azad Jammu \& Kashmir. The
system is not a model in a notebook. It ingests fresh observations every hour,
retrains itself every night, refuses to deploy a new model that is not measurably
better than the one already serving, explains every forecast it makes, watches
its own inputs for distributional drift, scores its own past forecasts against
the air quality that actually arrived, and serves all of this through a public
web dashboard and a documented programming interface. It does all of this on free
scheduled infrastructure, at an ongoing compute cost of effectively zero.
Figure~\ref{fig:arch-intro} shows the system at a glance.

\fig{architecture}{1.0}{High-level architecture of the \projname: three decoupled
pipelines on a scheduled compute plane, communicating only through a managed
feature store and model registry, with the results served by a container-hosted
API and a static dashboard.}{fig:arch-intro}

Figure~\ref{fig:arch-intro} divides the system into three regions. On the left, a
single free and keyless data provider supplies both pollutant and weather
observations, and also the weather \emph{forecast} that the inference stage needs.
In the centre, a compute plane of scheduled, disposable jobs runs the three
pipelines that give the architecture its name. On the right, a managed feature
store and model registry hold every piece of durable state, so that no pipeline
depends on any other pipeline still being alive. Beneath them, an application
programming interface and a browser application deliver the result to the public.
The remainder of this thesis develops each of these components in detail.

\section{Relevance to Course Modules}
The project is deliberately broad and exercises the core of a computer science
and data science curriculum. Table~\ref{tab:relevance} maps the principal
concepts the project applies to the areas in which they are taught.

\begin{xltabular}{\textwidth}{Q{3.6cm} Y}
\caption{Relevance of the project to the principal subject areas.}\label{tab:relevance}\\
\toprule
\bfseries Module / Area & \bfseries Concepts applied in this project \\
\midrule
\endfirsthead
\toprule
\bfseries Module / Area & \bfseries Concepts applied in this project \\
\midrule
\endhead
\bottomrule
\endlastfoot
Machine Learning & Supervised regression, gradient-boosted trees, regularised
linear models, bagged ensembles, hyper-parameter choice, bias-variance
reasoning, baseline comparison. \\
Time-Series Analysis & Autocorrelation, lag and rolling features, seasonality,
cyclical encodings, direct multi-horizon forecasting, chronological validation,
walk-forward backtesting, persistence baselines. \\
Deep Learning & A recurrent (LSTM) sequence model trained and evaluated on the
same task, and the reasoned decision not to promote it. \\
Data Engineering & Ingestion from a public API with retry and backoff, chunked
historical backfill, idempotent upserts, primary keys and event time, a feature
store and feature views, and a backend-agnostic storage facade. \\
Software Engineering & Layered package design, configuration as a single source
of truth, dependency inversion at the storage boundary, unit testing, iterative
lifecycle, and documentation. \\
MLOps and DevOps & Scheduled pipelines, continuous integration, a model registry,
champion-challenger promotion, artefact versioning, containerisation, and
production deployment. \\
Statistics & The Population Stability Index, empirical residual quantiles and
split-conformal prediction intervals, correlation analysis, distributional
comparison. \\
Explainable AI & Shapley additive explanations, global and per-instance
attribution, and the translation of attributions into natural language. \\
Web Technologies & A typed asynchronous HTTP API, cross-origin policy, a
single-page application, responsive charting, and interactive cartography. \\
Human-Computer Interaction & Information hierarchy on a public dashboard,
colour semantics for a standardised health scale, uncertainty communication, and
graceful degradation when a capability is unavailable. \\
\end{xltabular}

\section{Project Background}
Air-quality forecasting has three broad traditions. The first is
\emph{deterministic chemical transport modelling}: physically simulating the
emission, transport, chemistry and deposition of pollutants over a gridded
domain. These models are the scientific gold standard and underpin national
forecasting services, but they require an emissions inventory, meteorological
boundary conditions and considerable computing capacity, none of which is
available to an individual developer. The second is \emph{classical statistical
time-series modelling}, in which an autoregressive or moving-average model is
fitted to the pollutant series for a single site. Such models are cheap and
interpretable, but they capture the persistence of the series rather than its
drivers, and they degrade quickly beyond a short lead time because they cannot
use a weather forecast. The third, and the one this project belongs to, is
\emph{supervised machine learning}, in which a model learns the mapping from
weather, calendar and recent pollution state to future pollutant concentration
directly from historical data.

The machine-learning approach has one decisive practical advantage for the
problem at hand: the most informative predictors of tomorrow's air quality -
tomorrow's wind, temperature, humidity and rainfall - are themselves available
as a free public forecast. Air quality is, to a first approximation, a story
about dispersion. Wind clears pollutants; a temperature inversion traps them
under a warm lid; rain scavenges particulates out of the air column. If the
weather for the next three days can be obtained at no cost, then a model that has
learned how weather and pollution interact can produce a genuinely forward-looking
forecast without any physical simulation at all.

What the machine-learning approach does not solve by itself is the problem of
\emph{staying} correct. A model trained on last year's data and left alone will
decay, because the relationship it encodes is not stationary: seasons turn,
emission patterns change, and the distribution of the inputs moves away from the
distribution the model was fitted on. A great many student and research projects
stop at a trained model and a reported accuracy figure. The interesting
engineering problem - and the one this project treats as its centre of gravity
- begins immediately afterwards. How does a model get retrained without a human
present? How does one prevent a nightly retrain from silently shipping a worse
model? How is the reigning model stored so that it survives the machine that
trained it? How does the system notice that the world has changed? How does a
user know why the system believes what it says? These are the questions of
machine-learning operations, or MLOps, and they are what this thesis is
principally about.

\section{Motivation}
Three observations motivated this project.

The first is a matter of public health. Twenty-two cities, spanning every
administrative unit of the country, are covered by this system; between them they
are home to a very large share of Pakistan's urban population. A forecast with
three days of lead time changes what a household can do with the information. It
converts a warning into a plan. The exploratory analysis presented in
Chapter~\ref{ch:design} shows exactly how large the effect being forecast is: the
mean AQI across the whole record sits above 140 in five of the 22 cities
studied, close to the threshold the EPA labels \emph{unhealthy}, and the
seasonal signal is pronounced: January averages 146.8 against 91.5 in April,
the cleanest month. This is not a marginal environmental nuisance; it is a recurring, seasonal,
predictable public-health event, and predictability is precisely what a
forecasting system can exploit.

The second is a matter of engineering discipline. There is a large gap between a
model that scores well on a held-out split and a system that keeps working
unattended. Closing that gap normally implies infrastructure and therefore cost,
which is why so much student work stops at the notebook. One goal of this project
was to demonstrate that the entire lifecycle - scheduled ingestion, a feature
store, automated retraining, a gated model registry, evaluation, drift monitoring
and public deployment - can be assembled from free tiers and scheduled runners,
and that doing so imposes real constraints which are themselves instructive. Some
of the most valuable findings in this thesis, reported in
Chapter~\ref{ch:impl}, came from failures induced by exactly those constraints:
an insert that hung because a free-tier executor stalled, a statistics profiler
that exhausted its memory on a 74-column table, and a deployment tool that
silently omitted the model file because the repository ignored it.

The third is a matter of trust. A forecast that cannot be interrogated is a
forecast that will not be believed, particularly when it disagrees with what a
person can see from their window. The system therefore never presents a bare
number. Each forecast carries an 80\% prediction interval, so the user can see
how confident the model is; each carries a plain-language explanation generated
from the model's own Shapley attributions, so the user can see what is driving
it; and the model's complete promotion history, per-horizon accuracy, backtest
results and current drift status are all published in the same interface, so the
claim of accuracy can be checked rather than merely asserted.

\section{Problem Statement}
Existing air-quality information for Pakistan is current-only, sparse,
unexplained and manually maintained, and the machine-learning work that addresses
it typically stops short of an operating system. Official and commercial sources
report the AQI as it stands at a monitoring station, which tells a user what has
already happened rather than what is about to; where a forecast is offered it
generally covers a single major city, so residents of the smaller cities in
Khyber Pakhtunkhwa, Balochistan, Gilgit-Baltistan and Azad Jammu \& Kashmir are
served by nothing at all. No mainstream source explains its number, so a user
cannot distinguish a reading driven by a genuine pollution episode from one
driven by an unusual meteorological configuration, and none publishes its own
error, so its reliability cannot be assessed.

The research literature closes part of this gap and leaves the rest open.
Supervised models for pollutant forecasting are well established, and
gradient-boosted trees in particular are repeatedly found competitive with, or
superior to, deep sequence models on tabular, weather-driven environmental data.
What is far less commonly addressed is the operational half of the problem. A
published accuracy figure is usually computed on a single held-out split and
reported as one average number, which conceals the fact that error grows steeply
with lead time and thereby overstates usefulness at the horizons a user actually
cares about. Retraining, if discussed at all, is assumed rather than implemented;
promotion is unguarded, so a scheduled retrain can degrade a service without
anyone noticing; the trained artefact is left on the machine that produced it,
which in a serverless setting means it is destroyed minutes later; and drift,
which is the specific mechanism by which such a system decays, is rarely
instrumented.

There is therefore a need for a system that forecasts AQI several days ahead for
many cities rather than one; that reports its accuracy honestly, per horizon and
under rolling validation rather than a single split; that retrains itself without
human intervention while refusing to promote a model that is not demonstrably
better; that stores its champion durably and independently of any single machine;
that quantifies its own uncertainty and watches its own inputs and outputs for
decay; that explains each individual prediction in language a non-specialist can
act on; and that delivers all of this to the public reliably and at negligible
cost. A fuller treatment of the problem, of the shortcomings of current systems
and of the specific gaps this project addresses is given in
Chapter~\ref{ch:problem}.

\section{Aims and Objectives}
The overarching aim of the project is to design, implement, deploy and evaluate a
serverless, self-monitoring machine-learning system that forecasts the Air
Quality Index three days ahead for 22 Pakistani cities and explains its
predictions. This aim is decomposed into the following objectives.

\begin{enumerate}
  \item \textbf{Data and features.} To ingest weather and pollutant observations
  automatically from a free public source, and to engineer from them a rich but
  strictly leakage-controlled feature representation that generalises across many
  cities rather than being fitted to one.
  \item \textbf{Serverless pipeline architecture.} To realise the system as three
  decoupled Feature, Training and Inference pipelines, orchestrated by a free
  scheduler and communicating only through a managed feature store and model
  registry, so that each may run, fail and be replaced independently.
  \item \textbf{Accurate multi-horizon forecasting.} To train a single global
  model that serves every lead time from one hour to 72 hours by treating the
  horizon as an input feature, to measure it against a persistence baseline at
  every horizon rather than only on average, and to attach calibrated prediction
  intervals to every forecast it produces.
  \item \textbf{MLOps discipline.} To implement a champion-challenger promotion
  gate backed by a persistent leaderboard and a durable versioned model registry,
  and to evaluate the model rigorously through per-horizon metrics and
  walk-forward backtesting rather than a single averaged score.
  \item \textbf{A trust layer.} To explain each forecast through Shapley
  attributions rendered as plain language, to allow a user to interrogate the
  model through a what-if simulator, and to monitor the system continuously for
  input drift and realised forecast error.
  \item \textbf{Delivery.} To serve the whole system publicly through a typed
  HTTP interface and an interactive dashboard, and to expose the same grounded
  data to a language-model advisor through a standard tool protocol.
\end{enumerate}

\section{Scope of the Project}
The project delivers a complete, running system. In scope are: hourly ingestion
and feature engineering for 22 cities; a full historical backfill from January
2023 to the present; exploratory analysis of the resulting dataset; a daily
training pipeline comparing four model families; a global multi-horizon
forecaster covering lead times from one to 72 hours; empirical prediction
intervals; a feature store and a model registry; champion-challenger promotion
with a persistent leaderboard; per-horizon evaluation and walk-forward
backtesting; Shapley-based explanation in natural language; a what-if simulator;
drift and forecast-error monitoring; hazard alerting; a FastAPI backend; a React
dashboard; a Model Context Protocol server and a language-model advisor; and live
public deployment of the whole.

The following are explicitly out of scope. The project does not train a bespoke
foundation model or perform physical chemical-transport simulation; it uses a
hosted language model through its API for the advisor and a free public data
service for its observations. It does not attempt street-level or sub-hourly
resolution: the finest granularity of the data source, and therefore of the
system, is one city and one hour. It does not guarantee accuracy during
unprecedented events - an uncontrolled industrial fire, for example - that have
no analogue in the training record. It does not offer a mobile application, and
it does not use any paid data source. Finally, the system forecasts the AQI; it
does not attribute pollution to sources, and the what-if simulator is
explicitly presented as a model simulation rather than a causal claim.

\section{Methodology and Software Lifecycle}
The project followed an iterative and incremental lifecycle, organised as a
sequence of modules, and illustrated in Figure~\ref{fig:sdlc}. Each iteration
delivered one independently runnable and independently testable capability, and
each was allowed to inform the plan for the next. This was not merely a
convenient way to organise the work; it was the only defensible one, because the
behaviour of several components could be established only empirically. Whether a
free-tier feature store could absorb a three-and-a-half-year backfill for 22
cities, whether a sequence model would beat a tree ensemble on this data, and
whether a nightly retrain would in fact produce an improvement rather than noise
were all questions that had to be answered by running the thing rather than by
reasoning about it in advance.

\fig{sdlc}{0.72}{The iterative, module-by-module lifecycle followed during
development. Each iteration delivers one working, testable increment, and each
observed failure feeds directly into the next iteration's plan.}{fig:sdlc}

\subsection{Planning}
Each iteration began by selecting the next module and identifying its principal
risk. Risk drove ordering: the historical backfill, which depended on an external
free-tier service absorbing hundreds of thousands of rows, was attempted early
precisely because a failure there would have invalidated the architecture.

\subsection{Analysis}
Requirements for the increment were expressed as concrete scenarios - ``a
resident of Multan wants to know whether Thursday afternoon is safe for a child's
football practice'' - and reduced to functional and non-functional requirements,
which Chapter~\ref{ch:requirements} sets out in full.

\subsection{Design}
Interfaces were fixed before code was written, with particular care at the two
boundaries that determine whether the architecture holds: the contract between a
pipeline and the feature store, and the contract between a training run and the
model registry. Chapter~\ref{ch:design} documents these.

\subsection{Implementation}
Each module was implemented as a small, separately importable Python module with
one responsibility, so that ingestion, feature construction, training,
evaluation, explanation, monitoring and serving could each evolve without
disturbing the others. Chapter~\ref{ch:impl} presents the implementation.

\subsection{Testing and Integration}
Each increment was covered by unit tests over its critical logic - most
importantly the tests that assert the absence of target leakage - and then run
end to end against live data on a scheduled runner before being considered
complete. Chapter~\ref{ch:testing} reports the results.

\subsection{Maintenance}
Between iterations, the system was hardened in response to observed failures.
The backfill gained non-blocking writes and a history-aware resume after it froze
partway through a run; the training pipeline gained an explicit datetime
coercion after the feature store returned its event-time column as strings; and
the deployment gained explicit negations in the ignore file after the deployment
tool silently omitted the model bundle.

\subsection{System Development Lifecycle Summary}
The iterative approach suited a project whose central components could only be
judged by measurement. Each shortcoming observed in a real run became the first
item of the next iteration's plan, which produced steady and demonstrable
improvement - most visibly in the model itself, whose first champion was
promoted at a validation RMSE of 20.60 and then displaced by a better one at
19.69 once the full historical record was available.

\section{The Feature-Training-Inference Pattern in Brief}
The architectural backbone of the system is the Feature-Training-Inference
pattern, shown in Figure~\ref{fig:fti-intro}. Rather than a single program that
loads data, trains a model and predicts, the work is divided into three pipelines
that never call one another. The feature pipeline turns raw observations into
features and writes them to a feature store. The training pipeline reads
features from that store, produces a model and writes it to a model registry. The
inference pipeline reads the model from that registry and produces forecasts.
Each runs on its own schedule; each can fail, be re-run or be rewritten without
the others being aware of it; and the store and the registry are the only
contract between them.

\fig{fti-pipelines}{1.0}{The three decoupled pipelines, their schedules and the
shared state through which they communicate.}{fig:fti-intro}

This decoupling is what makes a serverless implementation possible. Because the
only durable state lives in the feature store and the model registry, the compute
itself can be entirely disposable: every pipeline runs on an ephemeral runner
that is destroyed when the job finishes, and nothing of value is lost when it is.

\section{Significance of the Project}
The significance of this work is threefold. \emph{Practically}, it puts a
three-day, explained, uncertainty-quantified air-quality forecast into the hands
of residents of 22 Pakistani cities, most of which are served by no forecast at
all, and it does so at a running cost that is effectively zero - which matters,
because a public-interest service that is expensive to run is a service that
stops running. \emph{Technically}, it demonstrates that the full MLOps lifecycle
- feature store, scheduled retraining, gated promotion, durable registry,
rigorous evaluation and self-monitoring - can be implemented in its entirety on
free infrastructure, and it documents the concrete constraints that such
infrastructure imposes and how each was overcome. \emph{Academically}, it
provides a fully documented reference design for an operational forecasting
system, including the components that are most often omitted from published work:
the promotion gate, the per-horizon and walk-forward evaluation that reveals how
skill actually decays with lead time, and the drift and realised-error monitoring
that tells an operator when the system has stopped being trustworthy.

\section{Thesis Organisation}
The remainder of this thesis is organised as follows.
Chapter~\ref{ch:literature} reviews related work in air-quality forecasting,
feature stores and the Feature-Training-Inference pattern, MLOps practice,
explainable machine learning and time-series validation.
Chapter~\ref{ch:problem} defines the problem and the gaps in existing systems.
Chapter~\ref{ch:requirements} presents the requirement analysis.
Chapter~\ref{ch:design} details the system design and architecture, and
Chapter~\ref{ch:impl} its implementation. Chapter~\ref{ch:ui} describes the user
interface. Chapter~\ref{ch:testing} reports the testing and the empirical
evaluation. Chapter~\ref{ch:conclusion} concludes and outlines future work, and
Chapter~\ref{ch:glossary} collects the abbreviations and glossary, followed by
the references.

\section{Summary of Introduction}
This chapter introduced the \projname\ and motivated it against the current state
of air-quality information in Pakistan, which is a nowcast for a few cities
rather than a forecast for many. It stated the problem, set out the aims,
objectives and scope, described the iterative methodology by which the system was
built, and previewed the Feature-Training-Inference pattern that gives the
architecture its shape. The next chapter surveys the body of work on which the
project builds.
